% Aula 06 --- quem é sensível à escala das covariáveis, e por quê.
Um colega recebeu um banco de dados em que a coluna \texttt{renda} está em reais e a coluna
\texttt{altura} em metros, e pergunta se precisa padronizar antes de ajustar. A
resposta honesta é "depende do método" --- e o critério que decide cabe numa
pergunta: \textbf{o método soma coisas vindas de colunas diferentes?}

Se soma, as unidades precisam ser comparáveis, e trocar uma delas estraga a
comparação. Se não soma, a unidade é irrelevante.
\begin{itens}
  \item \pts{2,0 --- 0,4 cada} Para cada método abaixo, diga se multiplicar uma
        única coluna por $1000$ \textbf{muda} ou \textbf{não muda} a predição,
        justificando em uma linha onde está --- ou onde não está --- a soma sobre
        colunas.
        \begin{enumerate}[label=(\roman*),leftmargin=1.0cm,nosep]
          \item Mínimos quadrados ordinários.
          \item Regressão Ridge com $\lambda>0$ fixo.
          \item Árvore de regressão.
          \item Floresta aleatória.
          \item KNN com $k$ fixo.
        \end{enumerate}
  \item \pts{0,9} Entre os que mudam, qual você espera que seja o mais afetado, e
        por quê? Compare o \emph{modo} como a escala entra em cada um.
  \item \pts{1,1} Outro colega ajustou o Lasso a dados padronizados, com $\lambda$
        escolhido por validação cruzada. Em seguida multiplicou por $1000$ a coluna
        de maior efeito sobre $Y$, reajustou com o \emph{mesmo} $\lambda$, e o erro
        quadrático médio no conjunto de teste \emph{diminuiu}. Ele conclui que é
        melhor não padronizar. Explique por que a reescala pode ter ajudado e por
        que a conclusão não se sustenta.
\end{itens}

\begin{solucao}
Uma única mudança de variáveis resolve de uma vez os três métodos lineares. Seja
$c=1000$, e suponha que a coluna reescalada é a primeira: os dados novos são
$\widetilde{\mathbb{X}}=\mathbb{X}D$, com $D=\operatorname{diag}(c,1,\dots,1)$. Para um
vetor $\bm\gamma$ de coeficientes nos dados novos, escreva $\bbeta=D\bm\gamma$, isto é,
$\beta_1=c\,\gamma_1$ e $\beta_j=\gamma_j$ para $j\ge2$. Então
$\widetilde{\mathbb{X}}\bm\gamma=\mathbb{X}\bbeta$: é \emph{o mesmo ajuste}, escrito em
outra unidade, e o RSS não muda. Já as penalidades mudam:
\[
  \sum_j\gamma_j^2=\frac{\beta_1^2}{c^2}+\sum_{j\ge2}\beta_j^2 ,
  \qquad
  \sum_j\abs{\gamma_j}=\frac{\abs{\beta_1}}{c}+\sum_{j\ge2}\abs{\beta_j}.
\]
Ou seja: \textbf{reescalar a coluna $1$ equivale a trocar o peso dela na penalidade}
--- por $1/c^2=10^{-6}$ no Ridge, por $1/c=10^{-3}$ no Lasso --- e deixar todo o resto
igual. Onde não há penalidade, não há o que mudar.

\smallskip
\textbf{(a)} Só dois mudam: o Ridge e o KNN.
\begin{enumerate}[label=(\roman*),leftmargin=1.0cm]
  \item \textbf{MQO: não muda.} Sem penalidade, o problema reescalado é o mesmo
        problema: o coeficiente da coluna vira $\beta_1/1000$, o produto $\beta_1x_1$
        fica igual, e as predições também. A predição \emph{soma} sobre as colunas,
        $\sum_j\beta_jx_j$, mas cada parcela é um número na unidade de $Y$, que o
        coeficiente ajusta livremente; nada na função-objetivo compara grandezas de
        colunas diferentes com pesos fixos.
  \item \textbf{Ridge ($\lambda>0$): muda.} A penalidade $\lambda\sum_j\beta_j^2$ soma,
        com o mesmo peso, coeficientes de colunas diferentes, e cada coeficiente carrega
        a unidade da sua coluna. Pela conta acima, a coluna reescalada passa a ser
        penalizada com peso $\lambda\cdot10^{-6}$.
  \item \textbf{Árvore: não muda.} Ela só pergunta "$x_j\le t$?", e
        $\{x_1\le t\}=\{1000\,x_1\le1000\,t\}$: as partições que se consegue formar são as
        mesmas, com o mesmo RSS cada uma, e o corte escolhido é o mesmo, com o limiar
        multiplicado por $1000$.
  \item \textbf{Floresta aleatória: não muda.} Cada árvore é indiferente à escala, e os
        sorteios --- das amostras bootstrap e das covariáveis candidatas --- não olham os
        valores. Com os mesmos sorteios, cada árvore sai idêntica, e a média também.
  \item \textbf{KNN com $k$ fixo: muda.} A distância $\sum_j(x_{ij}-x_{\ell j})^2$ soma,
        com o mesmo peso, diferenças medidas em unidades diferentes; a parcela da coluna
        reescalada fica multiplicada por $10^6$.
\end{enumerate}

\smallskip
\textbf{(b)} O \textbf{KNN}. O tamanho do fator não distingue os dois --- nos dois ele
chega a $10^6$ ---; o que distingue é \emph{o que a escala controla} em cada um.

No Ridge, a escala controla \emph{quanto um coeficiente é penalizado}. Pela conta do
início, o novo ajuste é o de um Ridge em que a coluna reescalada está praticamente
isenta de encolhimento. Todas as colunas continuam sendo usadas, e a mudança é limitada:
o pior que pode acontecer é aquele coeficiente deixar de ser encolhido --- e a diferença
some quando $\lambda\to0$, porque aí o Ridge vira o MQO, que não muda.

No KNN, a escala controla \emph{quem entra na média}. A distância passa a ser
\[
  10^6\,(x_{i1}-x_{\ell1})^2+\sum_{j\ge2}(x_{ij}-x_{\ell j})^2 ,
\]
e basta que duas observações difiram um pouco em $x_1$ para essa parcela dominar todas
as outras. Os vizinhos passam a ser, na prática, os de um KNN que só enxerga a coluna
reescalada: as outras colunas não ficam só menos importantes --- saem da conta, e a
informação que carregavam se perde. Nenhum $k$ a recupera.

\smallskip
\textbf{(c)} \emph{Por que a reescala pode ter ajudado.} Pela conta do início, o Lasso
nos dados reescalados é o Lasso original com a penalidade da coluna reescalada
multiplicada por $10^{-3}$: justamente a coluna que mais importa ficou quase isenta de
encolhimento. E encolher um coeficiente grande é quase só viés. No caso ortonormal, em
que o Lasso tem fórmula fechada, um coeficiente com $\abs{\widehat\beta_j}$ bem acima de
$\lambda/2$ é deslocado de exatamente $\lambda/2$: um deslocamento fixo, que não reduz
variância nenhuma e acrescenta $\lambda^2/4$ ao erro quadrático daquele coeficiente.
Tirar esse viés pode baixar o erro de teste --- pode, não precisa.

\emph{Por que a conclusão não se sustenta.} O resultado dependeu de \emph{qual} coluna
foi reescalada. Faça o mesmo com uma coluna \emph{irrelevante} e o efeito tende a se
inverter: quase isenta da penalidade, ela deixa de ser zerada, recebe um coeficiente
ajustado ao ruído, e o erro de teste piora. Além disso, a comparação usou o mesmo
$\lambda$, escolhido por validação cruzada para os dados padronizados; nos dados
reescalados, nem é esse o $\lambda$ que a validação cruzada escolheria.

O ponto de fundo é que, nos dois casos, quem decidiu quanto regularizar cada coluna não
foi o analista: foi a \textbf{unidade de medida} em que o dado chegou. Um procedimento
cujo resultado depende de a coluna ter sido anotada em metros ou em milímetros não está
bem definido, e ter dado certo desta vez não muda isso. Se a conclusão é que aquela
coluna merece menos penalidade, essa é uma decisão a tomar com os dados --- por
validação cruzada ---, e não a deixar por conta da unidade.
\end{solucao}
